\chapter{Future Work}
\label{ch:futurework}

The central finding of this dissertation --- that an explicit coordination mechanism
buys essentially nothing over independent learners on a \grid{2} traffic network ---
is best read not as an endpoint but as a precisely located starting point. Two
directions follow directly, both placed beyond the scope of this submission by
compute constraints rather than by design.

\section{Network Scale-Up (\texorpdfstring{\grid{2} $\rightarrow$ \grid{5}}{2x2 to 5x5})}
The most direct test of the interpretation in Chapter~\ref{ch:discussion} is to
grow the network. If the coordination value is near zero because a \grid{2} grid is
small enough that incentives are already locally aligned, then enlarging the network
to \grid{5} (25 intersections) should widen the gap between MAPPO and IPPO: distant
intersections interact through longer traffic chains, non-stationarity intensifies,
and the centralized critic's global view should begin to pay off. This expectation
is consistent with the large-scale findings of Chu~et~al.~\cite{chu2020largescale}
and the network-level cooperation of CoLight~\cite{wei2019colight}. A \grid{5} SUMO
network skeleton was prototyped during this project but reverted; training it to
convergence at this scale (with a single GPU and per-run training times measured in
days) exceeded the available compute budget. Quantifying how coordination value
scales with network size is therefore the primary item of future work.

\section{Extending to a Social Dilemma}
Traffic occupies the pure-coordination end of the cooperation spectrum: incentives
are aligned, so there is nothing to defect over. The complementary experiment is to
move to the \emph{opposite} end --- a sequential social dilemma in which defection is
individually tempting but collectively harmful. A natural choice is an
apple-harvesting commons in the style of Leibo~et~al.~\cite{leibo2017ssd}, where
agents collect a slowly regrowing resource and over-harvesting permanently depletes
it. There, the reward-structure axis (individual versus shared rewards) that is
inert in traffic should become decisive, allowing a direct study of when shared
rewards are \emph{necessary} to sustain cooperation versus when independent learning
suffices. Contrasting the learned behaviours across the two environments would
complete the coordination--dilemma spectrum this project set out to chart.

\section{Efficiency--Equity Analysis}
Finally, the per-agent breakdowns collected here can support a fairness analysis
(e.g.\ a Gini coefficient over per-intersection outcomes) to determine whether high
system performance requires unequal treatment of individual intersections --- a
question that becomes more pointed at larger scales and in dilemma settings.
